\documentclass[a4paper,oneside,12pt]{extarticle}

\usepackage[utf8]{inputenc}
\usepackage[english]{babel}
\usepackage{amssymb,amsthm,mathtools,enumitem,geometry,hyperref,booktabs}
\usepackage{graphicx}
\usepackage{xcolor}
\usepackage{tikz}
\usetikzlibrary{calc,patterns,decorations.pathreplacing}
\usepackage{tikz-cd}
\usepackage{pgfplots}
\usepgfplotslibrary{fillbetween}
\pgfplotsset{compat=1.18}
\usepackage{float}

\geometry{lmargin=15mm,rmargin=15mm,tmargin=10mm,bmargin=10mm,includeheadfoot}
\pagestyle{plain}
\frenchspacing

\theoremstyle{plain}
\newtheorem{proposition}{Proposition}
\newtheorem{lemma}[proposition]{Lemma}

\theoremstyle{definition}
\newtheorem{definition}[proposition]{Definition}
\newtheorem{example}[proposition]{Example}

\theoremstyle{remark}
\newtheorem{remark}[proposition]{Remark}

\newcommand*{\E}{\mathbb{E}}
\newcommand*{\Var}{\mathrm{Var}}
\newcommand*{\mP}{\mathbb{P}}
\newcommand*{\mR}{\mathbb{R}}
\newcommand*{\mbF}{\mathcal{F}}
\newcommand*{\ve}{\varepsilon}


\begin{document}

\begin{titlepage}
\begin{center}

{\footnotesize\scshape\color{gray}
Working Notes in Quantitative Risk and Modeling
}

\vfill

{\LARGE\bfseries Homoscedasticity and Its Testing:\\[4pt]
From Conditional Expectation to the Breusch--Pagan Test}

\vspace{36pt}

{\normalsize Boris Garbuzov}

\vfill

{\Large\bfseries Part~1~/~12}

\vspace{6pt}

{\Large\bfseries Conditional Expectation and the Regression Function}

\vfill

{\small July 2026}

\end{center}
\end{titlepage}

\setcounter{tocdepth}{2}
\tableofcontents
\newpage

%--------------------------------------------------------------------
\section{Homoscedasticity in Elementary Regression Textbooks}
\label{sec:elementary_textbooks}
%--------------------------------------------------------------------

Three settings are usually conflated in introductory regression
courses: $Y_i$'s with fixed indices, the linear model with non-random
$X$, and the joint-distribution setup with random $X$.  All three
talk about ``homoscedasticity'', but they mean different things, and
the slippage between them is a frequent source of confusion.  This
chapter separates them cleanly before moving on to the measure-theoretic
formulation used in the rest of the text.

%--------------------------------------------------------------------
\subsection{Non-random $X$: homoscedasticity as equal variances}
\label{sec:nonrandom}
%--------------------------------------------------------------------

In introductory regression courses, the regressor $X$ is treated as
non-random: it is a fixed set of index values $x_1, \ldots, x_n \in \mR$.
For each index $x_i$ there is a target random variable $Y_i$, all defined
on the same probability space. We can write this as
\[
  Y_{x_1},\ Y_{x_2},\ \ldots,\ Y_{x_n},
\]
and if we drop the strange subscripts, we simply have a random vector
$(Y_1, \ldots, Y_n)$---$n$ random variables on a common space.

\begin{definition}
In the non-random-$X$ setting, the variables $Y_1, \ldots, Y_n$ are
called \textbf{homoscedastic} if
\[
  \Var(Y_1) = \Var(Y_2) = \cdots = \Var(Y_n).
\]
\end{definition}

This is a definition by flat: equal marginal variances. It does not yet
involve a model or errors.

\begin{remark}
This can be taken as the definition, but the Probabilist's view is that it is better
understood as a \emph{consequence} of a more primitive assumption about
the errors, as shown in Section~\ref{sec:linearnonrandom}.
\end{remark}

%--------------------------------------------------------------------
\subsection{The linear model with non-random $X$}
\label{sec:linearnonrandom}
%--------------------------------------------------------------------

The standard linear model is
\[
  Y_i = \beta_0 + \beta_1 x_i + \ve_i, \qquad i = 1,\ldots,n,
\]
where $x_1,\ldots,x_n$ are fixed numbers and $\ve_1,\ldots,\ve_n$ are
random errors. Here the $x_i$ are not random---they are parameters of
the model, not outcomes of a random variable.

\textbf{Homoscedasticity in this model} means
\[
  \Var(\ve_1) = \Var(\ve_2) = \cdots = \Var(\ve_n) = \sigma^2.
\]
This implies $\Var(Y_i) = \sigma^2$ for all $i$, since $\beta_0 + \beta_1 x_i$
is a constant.

\textbf{Independence structure.}
If $\ve_1,\ldots,\ve_n$ are independent, then $Y_1,\ldots,Y_n$ are also
independent. However, i.i.d.\ errors do \emph{not} imply i.i.d.\ responses:
\[
  \E Y_i = \beta_0 + \beta_1 x_i,
\]
so the $Y_i$ have different means whenever the $x_i$ differ. Independent
with unequal distributions is not i.i.d.

\begin{remark}
This contrast matters: in the non-random-$X$ setting, ``identically
distributed'' is simply not the right concept for the $Y_i$. The
meaningful homogeneity condition is constancy of \emph{variance}, not
of distribution.
\end{remark}

%--------------------------------------------------------------------
\subsection{Random $X$: the joint distribution model}
\label{sec:randomX}
%--------------------------------------------------------------------

Now suppose $(X, Y)$ is a random vector defined on a probability space
$(\Omega, \mbF, \mP)$. We observe an i.i.d.\ sample of pairs:
\[
  (X_1, Y_1),\ (X_2, Y_2),\ \ldots,\ (X_n, Y_n),
\]
where each pair $(X_i, Y_i)$ has the same joint distribution as $(X, Y)$,
and the pairs are mutually independent (in the full sense: the generated
$\sigma$-algebras $\sigma(X_i, Y_i)$ are independent).

\textbf{The i.i.d.\ paradox.}
In this model, $(X_1, Y_1) \overset{d}{=} (X_2, Y_2) \overset{d}{=} \cdots$,
and in particular $Y_1 \overset{d}{=} Y_2$. Yet in the non-random-$X$ model,
$Y_1$ and $Y_2$ have different distributions (different means). What is
going on?

The resolution is that these are genuinely different models. In the
random-$X$ model, each $(X_i, Y_i)$ is an independent draw from the
\emph{same} joint distribution; the $X_i$ vary randomly, and the
marginal distribution of $Y_i$ is the same for all $i$ by symmetry.
In the non-random-$X$ model, the $x_i$ are fixed constants; conditioning
on specific values $x_1 \neq x_2$ selects slices of the joint distribution
with different conditional means.

\begin{example}
Consider a single random variable $X$ with sample $(X_1, X_2)$
i.i.d.$\sim X$. Before observing, $X_1 \overset{d}{=} X_2$. Once we
observe the outcomes $(x_1, x_2)$, we might think of $X_1$ and $X_2$
as degenerate: $X_1 \sim \delta_{x_1}$, $X_2 \sim \delta_{x_2}$. If
$x_1 \neq x_2$ these are different distributions---no longer i.i.d.
The non-random-$X$ regression model is precisely this: we condition on
the observed $x_i$, and the i.i.d.\ structure of the pairs collapses.
\end{example}

This is not a paradox but a model choice. Fixing $X$ (conditioning)
and treating $X$ as random are different setups, with different
implications for what ``i.i.d.'' means.

\section{Conditional Expectation: Definition and Properties}
\label{sec:condexp}
%--------------------------------------------------------------------

\subsection{Definition}

Let $(X,Y)$ be a random vector on $(\Omega, \mbF, \mP)$ with $\E|Y| < \infty$.
A random variable $Z = Z(\omega)$ is called $\E[Y \mid X]$, the
\textbf{conditional expectation of $Y$ given $X$}, if:
\begin{enumerate}[label=(\arabic*)]
  \item $Z$ is $\sigma(X)$-measurable, and
  \item $\int_A Z\,d\mP = \int_A Y\,d\mP$ for all $A \in \sigma(X)$.
\end{enumerate}

\begin{remark}[Operator view]
The conditional expectation is a binary operator taking two random
variables $X$ and $Y$ as arguments and returning a new random variable
$Z = \E[Y \mid X]$. More precisely:
\[
  \E[\,\cdot \mid \cdot\,]: (X, Y) \;\longmapsto\; Z,
\]
where the output $Z$ is itself a random variable on $(\Omega, \mbF, \mP)$,
with values in $\mR$. The same operator, applied to the same pair
$(X,Y)$ but with a squared-residual argument, gives the conditional
variance (the Conditional Variance section (Part~6)).

When we pass to the Doob--Dynkin representation $Z = h(X)$, we are
factoring this operator through the values of $X$: the function
$h:\mR\to\mR$ encodes the same information as $Z$, but as a function
of the \emph{realised value} of $X$ rather than of $\omega$ directly.
\end{remark}

\subsection{Existence, uniqueness, and distribution}

\begin{proposition}[Existence]
Under $\E|Y| < \infty$, a random variable $Z$ satisfying (1)--(2) always
exists.
\end{proposition}

\begin{proof}[Sketch]
This is a consequence of the Radon--Nikodym theorem. Define a signed
measure $\nu(A) = \int_A Y\,d\mP$ on $\sigma(X)$. Since $Y$ is
integrable, $\nu$ is absolutely continuous with respect to $\mP|_{\sigma(X)}$.
By Radon--Nikodym, there exists a $\sigma(X)$-measurable density $Z$
such that $\nu(A) = \int_A Z\,d\mP$ for all $A \in \sigma(X)$. This $Z$
satisfies (1) and (2).
Reference: Rudin, \textit{Real and Complex Analysis}, Theorem~6.10;
Kallenberg, \textit{Foundations of Modern Probability}, Theorem~5.1.
\end{proof}

\begin{proposition}[Uniqueness up to null sets]
If $Z_1$ and $Z_2$ both satisfy (1)--(2), then $Z_1 = Z_2$ almost surely.
\end{proposition}

\begin{proof}
Both satisfy condition (2), so $\int_A Z_1\,d\mP = \int_A Z_2\,d\mP$
for all $A \in \sigma(X)$, hence $\int_A (Z_1 - Z_2)\,d\mP = 0$.
Set $A = \{Z_1 > Z_2\} \in \sigma(X)$. On $A$ the integrand $Z_1 - Z_2 > 0$,
so $\mP(A) = 0$. Similarly $\mP(Z_1 < Z_2) = 0$.
Hence $\mP(Z_1 \neq Z_2) = 0$.
\end{proof}

\begin{remark}[Distribution]
Since any two versions of $\E[Y \mid X]$ agree almost surely, they
induce the same pushforward measure on $\mR$. That is, all versions
have the same distribution. This follows immediately from uniqueness:
if $Z_1 = Z_2$ a.s., then $\mP(Z_1 \in B) = \mP(Z_2 \in B)$ for
every Borel set $B$.
\end{remark}

\subsection{The equivalence theorem}

The following theorem gives an alternative characterisation of the
conditional expectation that follows directly from the definition,
without any use of the Doob--Dynkin factorisation.

\begin{proposition}[Equivalence theorem]
Let $Z$ be $\sigma(X)$-measurable with $\E|Z| < \infty$, and set
$u = Y - Z$. Then:
\[
  Z = \E[Y \mid X] \quad \Longleftrightarrow \quad \E[u \mid X] = 0.
\]
\end{proposition}

\begin{proof}
($\Rightarrow$) Assume $Z = \E[Y \mid X]$. Then:
\[
\E[u \mid X] = \E[Y - Z \mid X] = \E[Y \mid X] - \E[Z \mid X] = Z - Z = 0,
\]
where the last equality uses that $Z$ is $\sigma(X)$-measurable, so
$\E[Z \mid X] = Z$.

($\Leftarrow$) Assume $\E[u \mid X] = 0$. Rearrange: $Z = Y - u$.
Taking conditional expectations:
\[
\E[Z \mid X] = \E[Y - u \mid X] = \E[Y \mid X] - \E[u \mid X] = \E[Y \mid X].
\]
Since $Z$ is $\sigma(X)$-measurable, $\E[Z \mid X] = Z$, giving
$Z = \E[Y \mid X]$.
\end{proof}

\textbf{Decomposition.}
Taking $Z = \E[Y \mid X]$ and $u = Y - Z$, we have $Y = Z + u$ with
$\E[u \mid X] = 0$. This holds with no further assumptions beyond
$\E|Y| < \infty$.

\subsection{Doob--Dynkin factorisation and the regression function}

The equivalence theorem above characterises $\E[Y \mid X]$ purely
through the definition. The Doob--Dynkin lemma adds further structure:
since $Z = \E[Y \mid X]$ is $\sigma(X)$-measurable, there exists a
Borel-measurable function $h: \mR \to \mR$ such that
\[
  Z(\omega) = h(X(\omega)) \quad \text{for all } \omega \in \Omega.
\]
The function $h$ is the \textbf{regression function} of $Y$ on $X$:
$h(x) = \E[Y \mid X = x]$ for (almost) every $x$.
The decomposition now reads $Y = h(X) + u$.

\textbf{Degrees of freedom: $Z$ vs.\ $h$.}

\textit{Freedom of $Z$.}
By the uniqueness result (Section~4.2), any two versions $Z_1, Z_2$ of
$\E[Y \mid X]$ agree $\mP$-almost surely. The freedom of $Z$ is
exactly: modification on $\mP$-null subsets of $\Omega$.

\textit{Freedom of $h$.}
Any two Borel functions $h_1, h_2$ satisfying $h_1(X) = Z$ a.s.\ and
$h_2(X) = Z$ a.s.\ must satisfy $h_1(X) = h_2(X)$ $\mP$-a.s., which
means
\[
  \mP\bigl(h_1(X) \neq h_2(X)\bigr) = 0,
  \quad\text{i.e.}\quad
  P_X\bigl(\{x : h_1(x) \neq h_2(x)\}\bigr) = 0.
\]
So $h$ is unique $P_X$-almost everywhere on $\mR$: two versions may
differ only on a $P_X$-null Borel set. The freedom of $h$ is:
modification on $P_X$-null subsets of $\mR$.

\textit{Comparing the two freedoms.}
The relationship between the two measures is $P_X(B) = \mP(X \in B)$
for every Borel $B \subseteq \mR$. Therefore:
\begin{itemize}
  \item A $\mP$-null set $N \subseteq \Omega$ corresponds, via $X$,
        to modifying $h$ on $X(N) \subseteq \mR$, which is a
        $P_X$-null set. So every freedom of $Z$ translates into a
        freedom of $h$.
  \item Conversely, any Borel set $B \subseteq \mR$ with $P_X(B) = 0$
        gives a freedom of $h$ (modify $h$ on $B$ arbitrarily), and
        the corresponding set $\{X \in B\}$ has $\mP$-measure zero,
        so this change is invisible to $Z$ as well.
  \item However, $h$ also has freedom on Borel sets $B$ with
        $P_X(B) = 0$ that are \emph{not} in the image of any $\omega$
        under $X$---for instance, any open interval on which $P_X$
        assigns zero mass. For such $B$, modifying $h$ on $B$ cannot
        even be described in terms of $Z$: the event $\{X \in B\}$ has
        $\mP$-measure zero, so the modification has no counterpart on
        $\Omega$ at all.
\end{itemize}

This last point is the precise sense in which $h$ has \emph{strictly
more} freedom than $Z$. The extra freedom consists of modifications
on those $P_X$-null Borel sets in $\mR$ for which $\mP(X \in B) = 0$,
including sets that $X$ may reach on a $\mP$-null set of $\omega$'s.

\begin{remark}[Support and its complications]
One is tempted to describe the extra freedom of $h$ as ``freedom
outside the support of $P_X$.'' The support of $P_X$ on $\mR$ is
standardly defined as the smallest closed set of full $P_X$-measure,
or equivalently $\mR \setminus \bigcup\{U \text{ open} : P_X(U)=0\}$.
This is well-defined on $\mR$ with the usual topology, because an
arbitrary union of open $P_X$-null sets is again open and has
$P_X$-measure zero.

However, using the support to delineate freedom is imprecise in two
ways. First, the support can contain $P_X$-null subsets (individual
points under a continuous distribution), and $h$ is free on those too.
Second, range and support are different: $X(\Omega)$ and
$\mathrm{supp}(P_X)$ need not coincide---$X$ could take values outside
the support on a $\mP$-null set of $\omega$'s, or conversely the
support could contain points that $X$ never reaches. The cleanest
statement is: $h$ is determined $P_X$-almost everywhere on $\mR$, and
any Borel set of $P_X$-measure zero is a set of freedom for $h$.
\end{remark}

\textbf{The linear regression assumption.}
The standard linear model imposes $h(x) = \beta_0 + \beta_1 x$ on all
of $\mR$. This fixes $h$ on $P_X$-null sets where it is otherwise
undetermined---an extrapolation assumption that the data cannot test.

%--------------------------------------------------------------------
\section{Measure-Dependent Quantities}
\label{sec:measure_dependent}
%--------------------------------------------------------------------

This chapter collects the conceptual foundations that underlie the
likelihood-based tests discussed later: the notion of a
measure-dependent function, the distinction between a statistic and
a random quantity, the push-forward construction, and a concrete
illustration of iterating measure-dependent maps.


This section develops the conceptual framework needed to make the
prediction story from the Consequences section (Part~6) precise. It
clarifies the distinction between a statistic, a random quantity, and
a measure-dependent function --- a distinction that is usually left
implicit in textbooks.

\subsection{The problem with $\E[Y_2 \mid Y_1]$ as a predictor}

From the i.i.d.\ forecast notes: given a single observation $Y_1$,
the MSE-optimal predictor of $Y_2 \overset{d}{=} Y_1$ is
$f^*(Y_1) = \E[Y_2 \mid Y_1]$. But wait --- can we actually use this
as a predictor?

The conditional expectation $\E[Y_2 \mid Y_1]$ depends on the joint
distribution $P$ of $(Y_1, Y_2)$. It is not the same function for all
measures. A predictor that requires knowing $P$ to evaluate is not
useful in practice, and is not a statistic in any standard sense.

\subsection{Random quantities and statistics}

\begin{definition}
Let $\{P_\theta : \theta \in \Theta\}$ be a parametric family of
distributions on $\mR^n$. A \textbf{random quantity} is a measurable
function
\[
  f: \mR^n \times \Theta \to \mR,
  \quad (x_1,\ldots,x_n,\theta) \mapsto f(x_1,\ldots,x_n,\theta).
\]
A random quantity $f$ is called a \textbf{statistic} (with respect to
this family) if it does not depend on $\theta$:
\[
  f(x_1,\ldots,x_n,\theta) = f(x_1,\ldots,x_n,\theta')
  \quad \text{for all } \theta,\theta' \in \Theta.
\]
\end{definition}

\begin{remark}[Two views on statistics]
The Probabilist prefers the textbook convention: a statistic is simply a
measurable function of the data, with no $\theta$ argument at all.
The definition above is equivalent when the family is large enough ---
if the family contains all possible distributions, then any function
that happens to not depend on $\theta$ trivially satisfies the
definition. The explicit $\theta$ argument is useful when we want to
\emph{distinguish} statistics from non-statistics within a specific
model family.
\end{remark}

\begin{remark}[Monotonicity with respect to the family]
The statistic property is monotone with respect to inclusion of
measure families.
\begin{itemize}
  \item \textbf{Downward monotonicity.} If $f$ is a statistic with
        respect to a family $\mathcal{P}$, then it is also a statistic
        with respect to every subfamily $\mathcal{Q} \subseteq
        \mathcal{P}$. A smaller family imposes fewer constraints on
        measure-independence, so any function that does not depend on
        the measure within $\mathcal{P}$ certainly does not depend on
        it within $\mathcal{Q}$.
  \item \textbf{Upward monotonicity.} If $f$ is \emph{not} a statistic
        with respect to $\mathcal{P}$ (i.e.\ it takes different values
        for some $P, P' \in \mathcal{P}$), then it is also not a
        statistic with respect to any larger family $\mathcal{Q}
        \supseteq \mathcal{P}$, since $\mathcal{Q}$ still contains
        the distinguishing pair $P, P'$.
\end{itemize}
In particular: if $f$ \emph{is} a statistic with respect to the full
family of all probability measures, it is a statistic with respect to
every subfamily (downward monotonicity applied to the largest possible
family). Conversely, if $f$ is \emph{not} a statistic with respect to
some family $\mathcal{P}$, it fails to be a statistic with respect to
every superfamily $\mathcal{Q}\supseteq\mathcal{P}$. Being a statistic
with respect to a singleton family $\{P\}$ is vacuous and tells us
nothing about larger families.
\end{remark}

\textbf{Examples.}
\begin{itemize}
  \item $f(Y_1) = Y_1$: a statistic. It does not depend on $P$.
  \item $f_P(Y_1) = \E_P[Y_2 \mid Y_1]$: not a statistic. It depends
        on the joint distribution $P$ of $(Y_1,Y_2)$ and is different
        for different $P$.
  \item $h_P(X) = \E_P[Y \mid X]$ (regression function): not a
        statistic. It depends on the joint distribution of $(X,Y)$.
        To be precise, we should write $h_P$ or $h_{P_{XY}}$.
  \item $\hat\sigma^2 = \frac{1}{n-2}\sum\hat u_i^2$: a statistic.
        Computable from data alone.
\end{itemize}

\subsection{Mixed random quantities and approximation}

Some objects depend on both the sample and the parameter --- they are
neither pure statistics nor pure parameters. Examples:
\begin{itemize}
  \item $\E_P[Y_2 \mid Y_1]$: depends on $Y_1$ (data) and on $P$
        (measure).
  \item $\Var(\hat\beta_1 \mid \mathbf{X}) = \sigma^2/S_{XX}$: depends
        on $\mathbf{X}$ (data) and on $\sigma^2$ (parameter).
  \item $h_P(X_i) = \E_P[Y_i \mid X_i]$: depends on $X_i$ (data) and
        on $P$ (measure).
\end{itemize}

\textbf{How do we approximate mixed quantities, and by what?}
In principle, any random quantity --- statistic or not --- can serve
as an approximation. But in practice, only statistics are usable: a
non-statistic requires knowing the measure to evaluate, which defeats
the purpose. The plug-in principle is the standard way to convert a
non-statistic into a statistic: replace the unknown measure-dependent
object by an estimate.

\textbf{On the best estimator of $\E_P Y$.}
Suppose we want to estimate the parameter $\theta_P = \E_P Y$ by a
function $g$ of data $X_1,\ldots,X_n$. If we require $g$ to be a
true statistic (measure-independent), then the MSE depends on $P$
and we must specify for which $P$ we optimise.

If instead we allow $g$ to depend on $P$, the optimum is trivial:
set $g_P \equiv \theta_P$ (a constant equal to the parameter). This
is formally a valid ``estimator'' --- a constant is a function of the
sample --- but it is measure-dependent and thus not a statistic.
It uses the parameter to estimate itself, which is circular.

This shows why requiring measure-independence (the statistic property)
is not merely a convention: it prevents the degenerate solution of
``estimating $\theta_P$ by $\theta_P$.''

\begin{remark}[Singleton family]
If the family of measures is a singleton $\{P\}$, then every random
quantity is trivially a statistic. In this case, any random quantity is a
statistic. Is that right? Yes --- and this also means the constant
$\theta_P$ is a valid statistic in that case. The degenerate solution
is only excluded when the family is non-trivial.
\end{remark}

\begin{remark}[Degenerate case]
If the family of measures is a singleton $\{P\}$, then every random
quantity is trivially a statistic: there is no $\theta$ to depend on.
This is consistent with the definition but loses the distinction between
parameters, statistics, and measure-dependent quantities.
\end{remark}

\subsection{MSE decomposition: the general case}

Let $P$ be a joint distribution of $(Y_1, Y_2)$ with $\E_P[Y_2^2]<\infty$.
For any random quantity $f$ (possibly measure-dependent), the prediction
MSE decomposes as follows.

Add and subtract the measure-specific optimal predictor
$f_P^*(Y_1) = \E_P[Y_2 \mid Y_1]$:
\begin{align*}
\E_P[(f(Y_1)-Y_2)^2]
&= \E_P\bigl[(f(Y_1)-f_P^*(Y_1)) + (f_P^*(Y_1)-Y_2)\bigr]^2 \\
&= \E_P[(f(Y_1)-f_P^*(Y_1))^2]
  + 2\,\E_P[(f(Y_1)-f_P^*(Y_1))(f_P^*(Y_1)-Y_2)] \\
&\quad + \E_P[(f_P^*(Y_1)-Y_2)^2].
\end{align*}

\textbf{The cross-term vanishes.}
\begin{align*}
&\E_P\bigl[(f(Y_1)-f_P^*(Y_1))\cdot(f_P^*(Y_1)-Y_2)\bigr] \\
&= \E_P\bigl[\E_P[(f(Y_1)-f_P^*(Y_1))\cdot(f_P^*(Y_1)-Y_2)\mid Y_1]\bigr] \\
&= \E_P\bigl[(f(Y_1)-f_P^*(Y_1))\cdot\E_P[f_P^*(Y_1)-Y_2 \mid Y_1]\bigr] \\
&= \E_P\bigl[(f(Y_1)-f_P^*(Y_1))\cdot 0\bigr] = 0,
\end{align*}
where we used $\E_P[f_P^*(Y_1)-Y_2 \mid Y_1]
= f_P^*(Y_1) - \E_P[Y_2\mid Y_1] = 0$.

\textbf{Conclusion.}
\begin{equation}
\label{eq:mse_decomp}
\E_P[(f(Y_1)-Y_2)^2]
= \underbrace{\E_P[(f(Y_1)-f_P^*(Y_1))^2]}_{\geq\,0}
+ \underbrace{\E_P[(f_P^*(Y_1)-Y_2)^2]}_{\text{irreducible floor}}.
\end{equation}
Denoting the measure-specific minimizer $f_P^*(Y_1) = \E_P[Y_2 \mid Y_1]$:
\[
  \min_f \E_P[(f(Y_1)-Y_2)^2] = \E_P[(f_P^*(Y_1)-Y_2)^2],
\]
and $f_P^*$ is a minimizer. But this minimizer depends on $P$ ---
different measures give different optimal predictors.

\subsection{The maxi-min problem}

Since the optimal predictor $f_P^*$ changes with $P$, a natural
question is: can we find a single statistic $f$ (measure-independent)
that performs well uniformly across all measures?

One way to formalise this is the \textbf{maxi-min criterion}: find the
measure $P$ that is hardest to predict, and evaluate the best
measure-specific predictor there:
\[
  \max_{P_{Y_1,Y_2}} \min_{f_P} \E_P\bigl[(f_P(Y_1) - Y_2)^2\bigr].
\]
This is the worst-case irreducible floor over all joint distributions.
It does not directly give a statistic, because both the $\max$ and the
$\min$ are taken over measure-dependent objects.

\begin{remark}
This criterion does not appear in standard textbooks on prediction or
regression. The usual approach is to fix a family of measures (e.g.\
i.i.d., or parametric time-series) and work within that family.
The maxi-min formulation is a minimax decision-theoretic perspective
that treats the choice of measure as an adversary's move.
\end{remark}

A more tractable variant restricts to a subfamily
$\mathcal{P} \subseteq \{P_{Y_1,Y_2}\}$:
\[
  \max_{P \in \mathcal{P}} \min_{f_P} \E_P\bigl[(f_P(Y_1) - Y_2)^2\bigr].
\]
Natural subfamilies (in increasing order of restriction):
\begin{itemize}
  \item All joint distributions of $(Y_1,Y_2)$.
  \item Product measures: $Y_1$ and $Y_2$ independent.
  \item I.i.d.\ subfamily: $P = P_Y \otimes P_Y$.
  \item Parametric time-series models (e.g.\ AR(1)).
\end{itemize}
On the i.i.d.\ subfamily the irreducible floor is $\Var Y$, achieved
by $f_P^*(Y_1) = \E Y$ --- but this minimum depends on $P$ through
$\E Y$, so the maxi-min over i.i.d.\ measures is
$\sup_{P_Y} \Var_{P_Y} Y$, which is unbounded without moment
constraints.



Restrict to the i.i.d.\ subfamily: $P = P_Y \otimes P_Y$, so
$Y_1$ and $Y_2$ are independent with the same marginal distribution.
Under i.i.d., $\E_P[Y_2 \mid Y_1] = \E_P Y_2 = \E Y$ (a constant).

\textbf{Predictor 1: $f(Y_1) = Y_1$ (identity, a statistic).}
\begin{align*}
\E[(Y_1-Y_2)^2]
&= \E[Y_1^2 - 2Y_1 Y_2 + Y_2^2] \\
&= \E Y_1^2 - 2(\E Y)^2 + \E Y_2^2
  \quad\text{(independence)}\\
&= 2\E Y^2 - 2(\E Y)^2 = 2\Var Y.
\end{align*}
The identity function is a statistic: it does not depend on the
measure. MSE $= 2\Var Y$.

\textbf{Predictor 2: $f_P(Y_1) = \E_P[Y_2\mid Y_1] = \E Y$
(measure-dependent constant).}
\begin{align*}
\E[(\E Y - Y_2)^2] = \Var Y.
\end{align*}
This is the measure-specific optimal predictor. MSE $= \Var Y$,
exactly half of the identity predictor's MSE.

\textbf{What this illustrates.}
The optimal predictor $f_P^*$ is twice as good as the statistic $Y_1$
on the i.i.d.\ subfamily. But $f_P^*$ requires knowing $\E Y$, which
depends on $P$. In practice, $\E Y$ is unknown and must be estimated.
The plug-in predictor --- replace $\E Y$ by $\bar Y$ or $Y_1$ ---
recovers the statistic and pays the price: MSE doubles.

\begin{remark}
The conditional expectation $\E[Y_2 \mid Y_1]$ does not characterise
the full joint distribution of $(Y_1,Y_2)$: it is one operator of the
joint distribution, and many distributions share the same conditional
expectation (recall: conditional expectation is a bivariate operator).
For the purpose of optimal point prediction, knowing $\E[Y_2 \mid Y_1]$
is sufficient --- we do not need the full measure.
\end{remark}

\subsection{Connection to regression}

The same structure appears in regression. The conditional variance
$\Var(\hat\beta_1 \mid \mathbf{X}) = \sigma^2/S_{XX}$ is a random
variable (under homoscedasticity) with a single unknown ingredient:
the parameter $\sigma^2$. The function $g = \sigma^2$ is known only
up to the parameter --- it is the analogue of $\E Y$ in the prediction
example.

To predict $\Var(\hat\beta_1 \mid \mathbf{X})$, we estimate $\sigma^2$
by $\hat\sigma^2$ (plug-in) and substitute. This introduces the same
gap as replacing $\E Y$ by $Y_1$: the predictor is now a statistic,
but it pays an additional plug-in cost beyond the irreducible floor.

%--------------------------------------------------------------------


\end{document}
